% SPDX-License-Identifier: Apache-2.0
% covers: Remote, Ask, Answer
\datasheet{Remote}{A peripheral whose behaviour is a program somewhere else: every transaction leaves on a channel and the answer comes back on another.}

\dsfacts{\code{lib/parts/src/remote.rs}}{\code{txhdl\_parts::remote::Remote}}%
{\code{//docs:examples}}

\subsection*{Function}

\code{Remote} sits on AXI-Lite like any other peripheral and answers
nothing itself. Every transaction it accepts leaves on \code{out} as
an \code{Ask}, and the answer comes back on \code{back} as an
\code{Answer}. What is at the other end is not the peripheral's
business: in a simulation it is a model in the same program, and on a
board it is a frame on a wire and a program on another machine.

That is what it is for. A peripheral being designed can be written as
software first, against the bus it will really sit on and with the
rest of the system running as it really runs; when the software is
right the hardware replaces it and nothing else changes. Issue 297
says what that is worth.

What the bus sees is a slow device, and the part is honest about it
rather than hiding it: one transaction is outstanding at a time, and
the channel carrying it is held until the answer arrives, which is
what AXI-Lite allows and what a program at the far end of a wire
needs.

\subsection*{Parameters}

\begin{center}\footnotesize
\begin{tabular}{@{}lp{0.7\textwidth}@{}}
\toprule
Parameter & Meaning \\
\midrule
\code{T} & How many cycles a transaction may go unanswered before the
peripheral answers the bus \code{SLVERR} itself. A device that never
answers must not be able to stop the bus. The tables use 40.
\code{ex\_remote} uses 1000, since a frame each way through two MACs
is about seven hundred cycles, and \code{Board} uses two million,
20~ms at 100~MHz, which is a round trip to another machine. \\
\bottomrule
\end{tabular}
\end{center}

\subsection*{Register map}

None. Every address in the peripheral's range is the program's to
decide, and the address is carried out whole in the \code{Ask}.

\dstables{Remote}

\subsection*{Behaviour}

\begin{itemize}
\item One transaction is outstanding at a time, which \code{busy}
holds, and \code{writing} says which kind it is.
\item A write is taken before a read. Both are offered on channels of
their own, so a stream of reads could otherwise leave a write waiting
for ever.
\item Each transaction carries a \code{tag}, which counts up as each
leaves. An answer is taken from \code{back} whenever it is offered,
and counts only when its tag is the outstanding transaction's.
\item That is what makes the timeout safe. A transaction unanswered
for \code{T} cycles is answered \code{SLVERR} to the bus and finished
with; if its answer then arrives, its tag is no longer the one
outstanding and it is dropped rather than taken for the answer to
whatever came next.
\item The tag moves on when a transaction is finished with, not when
it leaves, so the late answer above is told apart from the next
transaction's.
\item A read is answered on \code{r} with the word the answer brought,
and a write on \code{b}; either carries \code{SLVERR} when the answer
said the transaction failed or the wait ran out.
\end{itemize}

\subsection*{Verification}

\begin{itemize}
\item Four tests in \code{lib/parts/src/remote.rs}, in
\code{//lib/parts:parts\_test}.
\code{a\_word\_written\_to\_the\_program\_is\_read\_back\_from\_it} is
the plain case.
\code{the\_program\_says\_which\_transactions\_fail} is an answer that
reports a failure.
\code{a\_slow\_program\_is\_waited\_for} holds the answer back and
finds the bus still served.
\code{a\_program\_that\_never\_answers\_does\_not\_stop\_the\_bus} is
the timeout, which is the test the part exists to pass.
\item \code{//lib/examples:ex\_remote} runs the whole chain, a
\code{RemoteLink} and two MACs with a wire between them, and the build
simulates this netlist against that run under nvc and Verilator:
\code{//docs:remote\_sim\_remote\_remote\_tb\_test} and
\code{//docs:remote\_vsim\_remote\_test}.
\item \code{//cpu/vreteno:board\_test} runs the core against it on the
board's design: a program writes a word to \code{0x3300} and reads it
back, with a program in the test answering the frames.
\end{itemize}

\subsection*{Used by}

\code{ex\_remote}, and \code{Board} at \code{0x3300}, where a
\code{RemoteLink<1>} carries it to the flagship's Ethernet port.
Issue 297 is what it is for: a peripheral and its software
co-developed, with the far end across a wire.

\subsection*{Limits}

AXI-Lite only, and one transaction outstanding, so the bus waits for
the far end at the speed the far end answers. It carries no transport
of its own: what moves an \code{Ask} to a program on another machine
is \code{RemoteLink} and a MAC, or anything else that speaks those two
channels. The timeout is counted in cycles rather than in time, so a
design that slows its clock waits proportionally longer.
